\subsection{Denominators and integrality properties}
\label{sub.denominators}


The universal denominator statement given in formula~\eqref{Dn} above was found empirically
on the basis of the extensive numerical data for the $4_1$, $5_2$ and the $(-2,3,7)$ pretzel
knots presented in the appendix to this paper. In this section, we prove it for the denominators
of the power series defined in terms of Gaussian-type integrals in~\cite{DG2}.
This proof only applies to $\s\in\calP\RED_K$, since there is no such integral representation
for the trivial representation, but the corresponding denominator statement is true here also
and can in fact be strengthened because the power series in that case come
from the Habiro ring, as explained at the end of this section.

\begin{Theorem}
\label{thm.denom}
For each knot $K$, representation $\s\in\calP_K$, and number $\a \in \BQ/\BZ$, we have
\[%\label{AD}
D_n  \tiA^{(K,\s)}_{\a,n} \in \calO_S\big[\z_\a,c^{-1}\big],
\]
where \smash{$\Phi_\a^{(K,\s)}(h)$} is as in~{\rm \cite{DG2}}, $D_n$ is as in~\eqref{Dn},
\smash{$\tiA^{(K,\s)}_{\a,n}$} is as in~\eqref{Phiprec2}, $c$ is the denominator of~$\a$,~$\calO$ is the ring of integers of $F_\s$ and $S$ is a finite set of primes
of $F_\s$ that depends on $K$ but not on $n$ or on~$\a$.
\end{Theorem}

The first few values of $D_n$ are given by
\begin{align*}%\label{Dnvalues}
&   1, \ 24, \ 1152, \ 414720, \ 39813120,  \ 6688604160, \ 4815794995200, \  115579079884800, \\
&  22191183337881600, \ 263631258054033408000, \ 88580102706155225088000, \\
&   27636992044320430227456000, \ 39797268543821419527536640000, \ \dots.
\end{align*}
Campbell Wheeler pointed out to us that the above sequence appears (with no proof)
to equal to the sequence~\texttt{A144618} of the online-encyclopedia of integer
sequences~\cite{OEIS}, the latter related to Stirling's formula with half-shift
$D_n=\text{den}(a_n)$, where
\[%\label{stirling}
z! \sim \sqrt{2\pi} (z+1/2)^{z+1/2} {\rm e}^{-z-1/2}
\sum_{n=0}^\infty \frac{a_n}{(z+1/2)^n},
\qquad z \to \infty .
\]

The numbers $D_n$ grow rapidly, for example
\[
D_{50} = 2^{197} 3^{72} 5^{19} 7^{11} 11^5 13^4 17^3 19^2 23^2 29^1
 31^1 37^1 41^1 43^1 47^1
\]
or $D_{50}/24^{50}50!=5^7 7^3 11^1 13^1 17^1$. \big(In general, $D_n/24^nn!$ is an
integer whose $n$th root tends to~\smash{$\prod_{p\ge5}p^{2/(p-1)(p-2)}=1.8592481285\dots$}.\big)
We give the first 50 values of~$D_n$ by tabulating the ratio $\delta_n=D_n/3D_{n-1}$
(after removing the power of~2, and omitting the values equal to~1):

\begin{center}
\small
\begin{tabular}{|ll|ll|ll|ll|ll|ll|}
\hline
$n$ & $\delta_n$ & $n$ & $\delta_n$ & $n$ & $\delta_n$ & $n$ & $\delta_n$ &
$n$ & $\delta_n$ & $n$ & $\delta_n$
\\ \hline
3 & $3 \tcdot 5$ &
11 & $13$ &
20 & $7$ &
27 & $3^3 \tcdot 5 \tcdot 11 \tcdot 29\!$ &
35 & $7^2 \tcdot 37$ &
42 & $3 \tcdot 5 \tcdot 23$
\\ %\hline
5 & $7$ &
12 & $3 \tcdot 5$ &
21 & $3 \tcdot 5 \tcdot 23\!$ &
29 & $31$ &
36 & $3^2 \tcdot 5 \tcdot 11$ &
44 & $13$
\\ %\hline
6 & $3 \tcdot 5$ &
15 & $3 \tcdot 5^2 \tcdot 7 \tcdot 17\!$ &
22 & $13$ &
30 & $3 \tcdot 5^2 \tcdot 7 \tcdot 17$ &
39 & $3 \tcdot 5 \tcdot 41$ &
45 & $3^2 \tcdot 5^2 \tcdot 7 \tcdot 11 \tcdot 17 \tcdot 47\!$
\\ %\hline
9 & $3^2 \tcdot 5 \tcdot 11\!$ &
17 & $19$ &
24 & $3 \tcdot 5$ &
33 & $3 \tcdot 5 \tcdot 13$ &
40 & $7$ &
48 & $3 \tcdot 5$
\\ %\hline
10 & $7$ &
18 & $3^2 \tcdot 5 \tcdot 11$ &
25 & $7$ &
34 & $19$ &
41 & $43$ &
50 & $7$
\\ \hline
\end{tabular}

\end{center}

\noindent
We also remark that $D_{n_1}D_{n_2}|D_{n_1+n_2}$ for all $n_1,n_2\ge0$ and hence
that the subgroup
\[%\label{defRD}
R_D[[h]] = \left\{ \sum_{n=0}^\infty \frac{a_n}{D_n} h^n  \,\big|\,
a_n \in R \  \text{for all} \   n \right\}
\]
of $K[[h]]$ is a subring for every subring $R$ of a field~$K$ of characteristic zero.

\begin{proof}
We give the proof only for the case~$\a=0$, $c=1$, using the formulas in~\cite{DG}.
The general case can be proved along the same lines using the more complicated
formulas in~\cite{DG2}, in which the Bernoulli numbers are replaced by Bernoulli
polynomials, but we do not give the details here.
We will also ignore the prime~2 in our proof since it behaves somewhat differently
and in any case can be added to the finite set of excluded primes~$S$ in the statement
of the theorem.

The power series \smash{$\Phi_0^{(K,\s)}(h)$} attached to a triangulation $\calT$
of~$\mathbb S^3\ssm K$ were defined in~\cite{DG} as formal Gaussian integrals
$\langle f_{\calT}\rangle$
of the formal power series
\be
\label{psixz}
f_{\calT}(x;z) = \exp\Bigg(\sum_{j=1}^N \sum_{\substack{r, k\ge0 \\ 2r+k-2 > 0}}
\frac{B_r}{r!} \frac{(-x_j)^k}{k!} \Li_{2-r-k}(z_j)  h^{r+\tfrac k2-1} \Bigg)
\ee
in a multi-variable $x=(x_1,\dots,x_N)$, where $z_1,\dots,z_N$ are the shape parameters
of~$\calT$ and where~${\langle f\rangle=\langle f(x)\rangle_Q}$ is the mean value
defined by Gaussian integration with respect to a certain quadratic form~$Q$ with
coefficients in the field~$F_\s$. This form is essentially the one given by the
symmetric matrix $A+B \,\diag(z_j/(1-z_j))$ that occurred in the discussion
of~\eqref{Nahm3}
in Section~\ref{sub.PhiPerturbative}, and the function~\eqref{psixz} is essentially
the product of the functions occurring in~\eqref{newpsixz}, except that the terms
with $r+k=m$ fixed were combined there into a single Bernoulli polynomial $B_m(x)$
for~$m\ge3$ or $B_2(x)-x^2$ for~$m=2$, and that the normalizations used in~\cite{DG}
were slightly different from the ones used in Section~\ref{sub.PhiPerturbative}.

We now expand the exponential in~\eqref{psixz} as the product of the exponentials
of the monomials in the sum, and recall that $\Li_{2-m}(z)\in\Z[1/(1-z)]$ for
every~$m\ge2$, to deduce that \smash{$\Phi_\a^{(K,\s)}(h)=\langle f_{\calT}\rangle$} is an
$R$-linear combination of Gaussian averages $\langle T\rangle$ of terms~$T$ of the
form
\be
\label{term}
T = \prod_{j=1}^N\prod_{\substack{r, k\ge0 \\ 2r+k\ge3\substack}}
\frac1{\l_j(r,k)!} \biggl(\frac{B_r}{r!}
\frac{x_j^k}{k!} h^{r+\tfrac k2-1}\biggr)^{\l_j(r,k)}
\ee
with non-negative multiplicities $\l_j(r,k)$, where $R=R_{\calT,\s}$ denotes the
ring generated over~$\Z$ by the numbers $(1-z_j)\i$. We write the monomial $T$ as
$c(T) h^n \bx^\bK/\bK!$ with
\be
\label{nAndK}
n = \sum_{j,r,k}\l_j(r,k)\left(r+\frac k2-1\right),\qquad K_j=\sum_{r,k}\l_j(r,k)k
\ee
and where for notational convenience have written $\bx$ and $\bK$ for the
$N$-tuples $(x_1,\dots,x_N)$ and~${(K_1,\dots,K_N)}$ (thus deviating from the
convention in the rest of the paper where boldface denotes matrices) and
$\bx^\bK/\bK!$ for the divided power \smash{$\prod x_j^{K_j}/K_j! $}. To prove the theorem,
we have to bound both the denominators of the numerical coefficient~$c(T)$ and the
further denominators coming from the Gaussian averaging
$ \bx^\bK/\bK! \mapsto\big\langle\bx^\bK/\bK! \big\rangle $.

We begin with the latter question. For this, we recall first that the Gaussian
average $\langle f\rangle_Q$ is given by \smash{${\rm e}^{\D_Q}(f)\bigl|_{x=0}$} for any power
series~$f$, where $\D_Q$ is the Laplacian associated to~$Q$, and hence is equal to
$\D_Q^\ell(f)/\ell!$ if $f$ is a homogeneous polynomial of degree $2\ell$. (For
polynomials of odd degree it of course vanishes trivially.) The Laplacian $\D_Q$
is a quadratic polynomial in the derivatives $\pt_i=\pt/\pt x_i$, and we can enlarge
the ring~$R$ by adjoining to it the coefficients of this polynomial, so since the
image of any divided factorial under any product $\pt_1^{\ell_1}\cdots\pt_N^{\ell_N}$
is an integer, we then certainly have that the Gaussian integral
$\big\langle\bx^\bK/\bK! \big\rangle$ is $1/\ell!$ times an element of~$R$. Unfortunately,
it turns out that this estimate is not good enough for our purposes, and we have
to work a little harder.

Writing $\D_Q$ as an $R$-linear combination of binomials $\pt_i\pt_j$ with
$1\le i\le j\le N$, and applying the multinomial theorem, we see that
$\D_Q^\ell/\ell!$ is an $R$-linear combination of terms
\smash{$\prod_{i\le j}(\pt_i\pt_j)^{m_{ij}}/m_{ij}!$} with $m_{ij}\in\Z_{\ge0}$. Define an even
symmetric $N\times N$ matrix $M$ by setting $M_{ij}=M_{ji}=m_{ij}$ for $i<j$ and
$M_{ii}=2 m_{ii}$. Then \smash{$\prod_{i\le j}(\pt_i\pt_j)^{m_{ij}}=\prod_j\pt_j^{K_j}$}
with $\bK=M\mb1$, where $\mb1$ is the vector consisting of~$N$~1's, and this
sends $\bx^\bK/\bK!$ to~1 and all other monomials to~0.
It follows that a universal denominator of $\big\langle\bx^\bK/\bK! \big\rangle$ is
the number
\begin{align*}
\D(\bK) := \text{l.c.m.}\Biggl\{ \prod_{1\le i\le N}(M_{ii}/2)!\prod_{1\le i<j\le N}M_{ij}!
 \,\Big|\, \text{$M=M^t\in M_{N,N}(\Z_{\ge0})$ even, $M\mb1=\bK$}\Biggr\} .
\end{align*}
Notice that this does divide~$\ell!$, because $\ell$ is the sum of the diagonal
$M_{ii}/2$ and of the $M_{ij}$ with~${i<j}$, so that this statement refines the bound
given above, but $\D(\bK)$ is in general much smaller than~$\ell!$ and this
improvement will be needed for the proof. A usually sharper multiplicative upper
bound for $\D(\bK)$ is the largest integer~$S$ whose square divides the product of
the~$K_j!$ (the proof of this also uses only the integrality of multinomial
coefficients), and then of course~$\D(\bK)$ also divides the g.c.d.\ of $\ell!$ and
$S$, which in general is smaller than either one. (For instance, for $\bK=(6,9,9,10)$
we have $\ell!=355687428096000$, $S=3135283200$, and~${(\ell!,S)=S/3}$.)
Either of these two latter upper bounds would be sufficient for our proof, but in
fact there is an easy upper bound that is stronger than either one of them and is
extremely sharp (in particular, it
is equal to $\D(\bK)$ for all $\bK$ with $N\le4$ and $\max(K_j)\le 30$), namely
\[
\D^*(\bK) := \prod_{\text{$p$ prime}} p^{\delta_p(\bK)} \quad \text{with} \quad
\delta_p(K_1,\dots,K_N) := \sum_{s\ge1}
\left[\frac12\sum_{j=1}^N\left[\frac{K_j}{p^s}\right]\right] .
\]
To show that $\D(\bK)$ divides~$\D^*(\bK)$, we need
$\sum_iV_p(M_{ii}/2)+\sum_{i<j}V_p(M_{ij})\le\delta_p(\bK)$ for every prime~$p$ and
every even symmetric matrix~$M$ with non-negative entries and row sums~$\bK$, where~${V_p(m):=v_p(m!)}$ for any~$m\ge0$ denotes the largest power of~$p$ dividing~$m!$.
In view of the standard formula $V_p(m)=\sum_{s\ge1}[m/p^s]$, it suffices for this
to show that
\[\sum_i[M_{ii}/2q]+\sum_{i<j}[M_{ij}/q]\le\frac12\sum_j[K_j/q]\]
for every prime power~$q$, and this follows immediately from the obvious facts
$[x/2q]=[[x/2]/q]$ and $[x]+[y]+\cdots \le[x+y+\cdots]$ valid for arbitrary real
numbers $x, y,\dots$.

Now going back to our main problem, we now see that it suffices to show that the
product~$\D^*(\bK) c(T)$ has denominator dividing~$D_n$ for all terms~$T$ as above,
with $\bK=(K_1,\dots,K_N)$ and~$n$ defined by~\eqref{nAndK}. We will prove this
one prime at a time (ignoring the prime~2), which is natural in view of the fact
that the upper bound~$\D^*(\bK)$ is defined by its prime power decomposition. To do
this, we will split both the term~$T$ and the corresponding numerical
coefficient~$c(T)$, and also each of the $N$ factors $T_j$ and $c(T_j)$ of which
they are comprised, as the product of four factors in a way depending on the
prime~$p$ being studied, labelled ``$s$'' (terms with $r=0$ and~$k$ smaller than~$p$),
``$b$'' (terms with $r=0$ and $k$ bigger than or equal to~$p$),
``1''~(terms with~${r=1}$), and ``$\ge2$'' (terms with $r\ge2$), with a similar splitting
of the individual weights~$K_j$ into the sum of four pieces
\smash{$K_j^{(s)}=\sum_{3\le k<p}\l_j(0,k) k$}, \smash{$K_j^{(b)}=\sum_{k\ge p}\l_j(0,k) k$}, \smash{$
K_j^{(1)}=\sum_{k\ge1}\l_j(1,k) k$}, and~\smash{$K_j^{(\ge2)}=\sum_{r\ge2, k\ge0}\l_j(r,k) k$}.
We also decompose the number~$n$ (the exponent of~$h$ in~$T$) in~\eqref{nAndK} as
$\ell+n'-t$ with
\[
\ell := \frac12 \sum_{j=1}^NK_j , \qquad
n' :=\sum_{\substack{1\le j\le N \\ r\ge2, k\ge0\substack}}\l_j(r,k) (r-1) , \qquad
t :=\sum_{\substack{1\le j\le N \\k\ge3\substack}}\l_j(0,k)
\]
and also split $t$ as \smash{$t^{(s)}+t^{(b)}$} according as $3\le k\le p-1$ or $k\ge p$ in
the summation, with each of \smash{$t^{(s)}$} and \smash{$t^{(b)}$} splitting into the sum over
$1\le j\le N$ of pieces \smash{$t_j^{(s)}$} and \smash{$t_j^{(b)}$} in the obvious way.

The numerical coefficient $c(T)$ can be decomposed as
\be
\label{cTdecomp}
c(T) = \prod_{1\le j\le N}\frac{K_j!}{P_j(T)} \cdot
\prod_{\substack{1\le j\le N \\ r\ge2, k\ge0\substack}}
\frac1{\l_j(r,k)!} \biggl(\frac{B_r}{r!}\biggr)^{\l_j(r,k)}
\ee
with
\[
P_j(T) = \prod_{0\le r\le1, k\ge1}\l_j(r,k)! k!^{\l_j(r,k)} \cdot
\prod_{r\ge2, k\ge0}k!^{\l_j(r,k)} ,
\]
which we can split up further as \smash{$P_j^{(s)}(T)P_j^{(b)}(T)P_j^{(1)}(T)P_j^{(\ge2)}(T)$}.
The reason that we have included the factor $\l_j(r,k)$ into the definition of
$P_j(T)$ for $0\le r\le1$ but not for~$r\ge2$ is that~$\l!k!^\l$ divides $(k\l)!$
for all $\l\ge0$ if $k$ is strictly positive but not if $k=0$, and the terms with~${r\ge2}$ can have~$k=0$. Then the product \smash{$P_j^{(b)}(T)P_j^{(1)}(T)P_j^{(\ge2)}(T)$}
divides \smash{$\big(K_j-K_j^{(s)}\big)!$}, while the first factor \smash{$P_j^{(s)}(T)$} divides \smash{$t_j^{(s)}!$}
up to a $p$-adic unit because here $k$ is always less than~$p$ and therefore~$k!$
is not divisible by~$p$. (Here we have made repeated use of the integrality of
multinomial coefficients.) On the other hand, by Lemma~\ref{LemmaB} below
and the submultiplicativity of~$D_n$, the second factor in~\eqref{cTdecomp} has
denominator dividing~$D_{n'}$. Putting this all together, we deduce that~$c(T)$
is $G(T)/D_{n'}$ times a $p$-adic integer for every~$p$ (always different from~2
and not dividing the denominators of the elements of~$R$), where
\smash{$G(T)=\prod_{j=1}^N\bigl(K_j!/t_j^{(s)}! \big(K_j-K_j^{(s)}\big)!\bigr)$}. Using the
submultiplicativity of $D_n$ again, this reduces the problem to showing that
${\delta_p(\bK)\le v_p(D_{\ell-t} G(T))}$ for each~$p$, and in view of
the definitions of $\delta_p(\bK)$ and $D_n$ and of the above-mentioned formula~${V_p(m)=\sum_{s\ge1}[m/p^s]}$ for the $p$-adic valuation of factorials, this in turn
will follow if we can show that
\be
\label{Final}
\Biggl[\frac12\sum_{j=1}^N\biggl[\frac{K_j}q\biggr]\Biggr] \le
\biggl[\frac{\ell-t}{q^*}\biggr]
+ \sum_{j=1}^N \Bigg(\biggl[\frac{K_j}q\biggr] \m \Bigg[\frac{t_j^{(s)}}q\Bigg]
\m \Bigg[\frac{K_j-K_j^{(s)}}q\Bigg] \Bigg)
\ee
for each prime power $q=p^s$ with $s\ge1$, where $q^*:=p^{s-1}(p-2)$. % =(1-2/p) q$.

For this final step, we first note that
\[
\frac{\ell-t}{q^*} = \sum_{j=1}^N
\frac{K_j^{(s)}+K_j^{(b)}+K_j^{(1)}+K_j^{(\ge2)}-2t_j^{(s)}-2t_j^{(b)}}{2q^*}
 \ge \sum_{j=1}^N\frac{K_j-2t_j^{(s)}}{2q}
\]
since
\[
\frac{\big(K_j^{(b)}-2t_j^{(b)}\big)}{q^*}\ge \left(1-\frac{2}{p}\right) \frac{K_j^{(b)}}{q^*}=\frac{K_j^{(b)}}{q}
\]
 \big(because $k\ge p$ in the terms defining \smash{$K_j^{(b)}$} and \smash{$t_j^{(b)}$}\big) and~$q^*<q$.
Using that $\big[\frac x{2q}\big]=\big[\frac12\big[\frac xq\big]\big]$, we deduce that
\[
\biggl[\frac{\ell-t}{q^*}\biggr] \ge
\Biggl[\frac12\sum_{j=1}^N\Bigg[\frac{K_j-2t_j^{(s)}}{q}\Bigg]\Biggr]
\]
and hence (since $x\le y$ certainly implies $[x/2]\le[y/2]$) the
inequality~\eqref{Final} will follow if we have the inequality
\[
\biggl[\frac{K_j}q\biggr] \le \Bigg[\frac{K_j-2t_j^{(s)}}{q}\Bigg]
+ 2 \Bigg(\biggl[\frac{K_j}q\biggr] \m \Bigg[\frac{t_j^{(s)}}q\Bigg]
\m \Bigg[\frac{K_j-K_j^{(s)}}q\Bigg] \Bigg)
\]
for every $1\le j\le N$. But this inequality is trivial, since
\smash{$K_j-K_j^{(s)}\le K_j-3t_j^{(s)}\le K_j-2t_j^{(s)}$} (because every $k$ in the
definition of \smash{$K_j^{(s)}$} is $\ge3$) and
\[
\biggl[\frac{K_j}q\biggr]\ge\Bigg[\frac{K_j-2t_j^{(s)}}q\Bigg]
+2\Bigg[\frac{t_j^{(s)}}q\Bigg].
\]
 This completes the proof of Theorem~\ref{thm.denom} modulo that of the following lemma.
\end{proof}

\begin{Lemma}\label{LemmaB}
For any integers $r\ge2$ and $\l\ge0$, we have
\be
\label{BernEstimate}
\frac1{\l!} \left(\frac{B_r}{r!}\right)^\l \in \frac1{D_{\l(r-1)}} \Z .
\ee
\end{Lemma}

\begin{proof}
We prove this one prime at time. By well-known results of
von Staudt and Clausen, the Bernoulli number $B_r$ ($r>0$) has $p$-adic valuation
$-1$ if~${(p-1)|r}$ and $B_r/r$ is $p$-integral if~${p-1}$ does not divide~$r$. From this
we deduce that the $p$-adic valuation of the denominator of $B_r/r!$ is bounded
above by $\bigl[\frac r{p-1}\bigr]$, which is $\le\bigl[\frac {r-1}{p-2}\bigr]$
since $\frac{r-1}{p-2}\ge\frac r{p-1}$ if $r\ge p-1$ and both expressions
vanish otherwise. The $p$-adic valuation of the denominator of
\smash{$\frac1{\l!}\bigl(\frac{B_r}{r!}\bigr)^\l$} is therefore bounded above by
\smash{$V_p(\l)+\l\bigl[\frac{r-1}{p-2}\bigr]$}. On the other hand, from the definition
of~$D_n$ we have $v_p(D_n)=m+V_p(m)=V_p(pm)$, where $m=\bigl[\frac n{p-2}\bigr] $.
We must therefore show that
\[
V_p(\l)+\l \biggl[\frac{r-1}{p-2}\biggr] \le
V_p\biggl(p \biggr[\frac{\l(r-1)}{p-2}\biggr]\biggl)
\]
for all $\l\ge0$ and $r\ge2$. For this, we make a case distinction: if $ 2\le r<p-1$, then
\[
\text{l.h.s.} = V_p(\l) = V_p\biggl(p \biggl[\frac\l p\biggr]\biggr) \le
V_p\biggl(p \biggl[\frac\l{p-2}\biggr]\biggr) \le \text{r.h.s.} ,
\]
while if $r\ge p-1$, then we set $h=\bigl[\frac{r-1}{p-2}\bigr]\ge1$ and have
instead
\[
\text{l.h.s.} = V_p(\l)+\l h \le \l h+V_p(\l h) = V_p(p\l h) \le \text{r.h.s.}
\]
because $[\l x]\ge \l [x]$ for any positive real number~$x$.
\end{proof}
